% Part: incompleteness
% Chapter: representability-in-q
% Section: prim-rec

\documentclass[../../../include/open-logic-section]{subfiles}

\begin{document}

\olfileid{inc}{req}{pri}
\olsection{आदिम पुनरावर्तनाचे अनुकरण}

बीटा फलन वापरून नियमित लघुतमीकरणाने आदिम पुनरावर्तनाने केलेल्या
व्याख्येचे ``अनुकरण'' करता येते, हे आता दाखवू. समजा आपल्याकडे
$f(\vec x)$ आणि $g(\vec x, y, z)$ ही फलने आहेत. मग $f$ आणि~$g$
यांपासून आदिम पुनरावर्तनाने परिभाषित केलेले फलन~$h(\vec x,y)$
असे आहे:
\begin{align*}
h(\vec x, 0) & =  f(\vec x) \\
h(\vec x, y+1) & =  g(\vec x, y, h(\vec x, y)).
\end{align*}
मूलभूत फलने आणि त्यांपासून संयोजन व नियमित लघुतमीकरणाने
परिभाषित केलेली फलने (उदाहरणार्थ~$\beta$) वापरून,
$f$ आणि~$g$ यांपासून केवळ संयोजन व नियमित लघुतमीकरणाने
$h$ परिभाषित करता येते, हे दाखवायचे आहे.

\begin{lem}
\ollabel{lem:prim-rec}
$h$ हे $f$ आणि $g$ यांपासून आदिम पुनरावर्तनाने परिभाषित
करता येत असेल, तर $f$, $g$ आणि $\Zero$, $\Succ$,
$\Proj{n}{i}$, $\Add$, $\Mult$, $\Char{=}$ ही फलने
वापरून संयोजन व नियमित लघुतमीकरणानेही ते परिभाषित करता येते.
\end{lem}

\begin{proof}
प्रथम $\hat h(\vec x, y)$ हे साहाय्यक फलन परिभाषित करू.
ते लघुतम संख्या~$d$ परत करते; हा $d$ पुढील अटी
पूर्ण करणाऱ्या क्रमिकेचा संकेतांक असतो:
\begin{enumerate}
\item $(d)_0 = f(\vec x)$, आणि
\item प्रत्येक $i < y$ साठी $(d)_{i+1} = g(\vec x, i, (d)_i)$,
\end{enumerate}
येथे $(d)_i$ हा आता $\beta(d,i)$ चा संक्षेप आहे. दुसऱ्या
शब्दांत, $\hat h$ हे $\tuple{h(\vec x, 0), h(\vec x, 1),
\dots, h(\vec x, y)}$ या क्रमिकेचा संकेतांक परत करते.
$\hat h$ पुढीलप्रमाणे लिहिता येते:
\[
\hat h(\vec x, y) = \umin{d}{(\beta(d,0) = f(\vec x) \land \bforall{i <
  y}{\beta(d,i+1) = g(\vec x, i,\beta(d,i)})}.
\]
लक्षात घ्या: येथे आदिम पुनरावर्तन लागत नाही; फक्त
लघुतमीकरण पुरेसे आहे. बीटा फलनावरील साहाय्यक विधान
\olref[bet]{lem:beta} नुसार या लघुतमीकरणातील शोध नियमित आहे.

आता आपल्याकडे
\[
h(\vec x, y) = \beta(\hat h(\vec x, y), y),
\]
म्हणून मूळ दिलेली दोन फलने आणि मूलभूत फलने यांपासून केवळ संयोजन
व नियमित लघुतमीकरणाने $h$ परिभाषित करता येते.
\end{proof}

\end{document}
